
\documentclass{article}
\usepackage{colortbl}
\usepackage{makecell}
\usepackage{multirow}
\usepackage{supertabular}

\begin{document}

\newcounter{utterance}

\centering \large Interaction Transcript for game `cladder', experiment `full\_v1.5\_default', episode 3456 with Qwen3.5{-}27B.
\vspace{24pt}

{ \footnotesize  \setcounter{utterance}{1}
\setlength{\tabcolsep}{0pt}
\begin{supertabular}{c@{$\;$}|p{.15\linewidth}@{}p{.15\linewidth}p{.15\linewidth}p{.15\linewidth}p{.15\linewidth}p{.15\linewidth}}
    \# & \multicolumn{2}{c}{Player} && \multicolumn{2}{c}{Game Master} \\
    \hline

    \theutterance \stepcounter{utterance}  
    & & & \multicolumn{4}{p{0.6\linewidth}}{
        \cellcolor[rgb]{0.9,0.9,0.9}{
            \makecell[{{p{\linewidth}}}]{
                \texttt{\tiny{[P1$\langle$GM]}}
                \texttt{You are an expert in causal inference. The following question is not a typical commonsense query, but rather a meticulously designed question created by a professor specializing in causal inference, intended to assess the students' mastery of the course content.} \\
\texttt{Imagine a self{-}contained, hypothetical world with only the following conditions, and without any unmentioned factors or causal relationships: Solar eclipse has a direct effect on Alice arriving to school. Traffic has a direct effect on Alice arriving to school. Method 1: We look at how solar eclipse correlates with Alice arriving to school case by case according to traffic. Method 2: We look directly at how solar eclipse correlates with Alice arriving to school in general. To understand how solar eclipse affects Alice arriving to school, is it more correct to use the Method 1 than Method 2?} \\
\texttt{Start your answer with 'Yes' or 'No', followed by additional reasoning or evidence to support your explanation.} \\
            }
        }
    }
     \\ \\

    \theutterance \stepcounter{utterance}  
    & \multicolumn{4}{p{0.6\linewidth}}{
        \cellcolor[rgb]{0.561,0.737,0.561}{
            \makecell[{{p{\linewidth}}}]{
                \texttt{\tiny{[P1$\rangle$GM]}}
                \texttt{No} \\
\\ 
\texttt{In the hypothetical world described, the causal structure is defined such that "Solar eclipse" (\$X\$) and "Traffic" (\$Z\$) are both direct causes of "Alice arriving to school" (\$Y\$), with no causal link between \$X\$ and \$Z\$ (since no unmentioned factors or relationships exist). This structure represents a **fork** or a **common effect** scenario where \$X\$ and \$Z\$ are independent causes of \$Y\$.} \\
\\ 
\texttt{Method 2 (looking at the general correlation) estimates the **total effect** of the solar eclipse on Alice's arrival. Since there is no confounding variable (a common cause of both the eclipse and traffic) and no mediator (traffic is not caused by the eclipse), the marginal association between the solar eclipse and Alice's arrival is an unbiased estimate of the causal effect.} \\
\\ 
\texttt{Method 1 (stratifying by traffic) attempts to control for \$Z\$. However, because \$Z\$ (Traffic) is a **collider** (it is a common effect of \$X\$ and potentially other unmentioned factors, but in this specific isolated world, it is simply an independent cause alongside \$X\$), conditioning on it can induce a spurious association if there were other causes, or more critically in this specific "no unmentioned factors" setup, it is simply unnecessary. More importantly, if we interpret the structure strictly as \$X \textbackslash{}to Y\$ and \$Z \textbackslash{}to Y\$ with \$X \textbackslash{}perp Z\$, conditioning on \$Z\$ does not introduce bias in the standard sense of confounding, but it reduces statistical efficiency (increases variance) without removing any bias.} \\
\\ 
\texttt{However, the question asks if Method 1 is *more correct*. In causal inference, "more correct" usually implies obtaining an unbiased estimate of the causal effect. Since there is no confounding, the marginal effect (Method 2) is the correct estimate of the total causal effect. Method 1 estimates the **direct effect** (or conditional effect) holding traffic constant. If the question implies understanding the *total* impact of a solar eclipse (which includes all pathways, though here there is only one direct path), Method 2 is the standard approach.} \\
\\ 
\texttt{Crucially, if the prompt implies a scenario where Traffic is a **collider** (e.g., if there were unobserved common causes, which the prompt explicitly forbids), conditioning would be bad. But even without unobserved causes, if \$X\$ and \$Z\$ are independent causes, conditioning on \$Z\$ (Method 1) yields the conditional effect \$E{[}Y|do(X), Z{]}\$, which is not necessarily the same as the total effect \$E{[}Y|do(X){]}\$ if there is an interaction between \$X\$ and \$Z\$. If there is no interaction, they are the same, but Method 1 is less efficient. If there *is* an interaction, Method} \\
            }
        }
    }
    & & \\ \\

    \theutterance \stepcounter{utterance}  
    & & & \multicolumn{2}{p{0.3\linewidth}}{
        \cellcolor[rgb]{0.9,0.9,0.9}{
            \makecell[{{p{\linewidth}}}]{
                \texttt{\tiny{[GM$|$GM]}}
                \texttt{No} \\
            }
        }
    }
    & & \\ \\

    \theutterance \stepcounter{utterance}  
    & & & \multicolumn{2}{p{0.3\linewidth}}{
        \cellcolor[rgb]{0.9,0.9,0.9}{
            \makecell[{{p{\linewidth}}}]{
                \texttt{\tiny{[GM$|$GM]}}
                \texttt{no} \\
            }
        }
    }
    & & \\ \\

    \theutterance \stepcounter{utterance}  
    & & & \multicolumn{2}{p{0.3\linewidth}}{
        \cellcolor[rgb]{0.9,0.9,0.9}{
            \makecell[{{p{\linewidth}}}]{
                \texttt{\tiny{[GM$|$GM]}}
                \texttt{game\_result = WIN} \\
            }
        }
    }
    & & \\ \\

\end{supertabular}
}

\end{document}
